\chapter{Pengujian dan Evaluasi Kebutuhan Sistem}
\label{lampiran:pengujian-dan-evaluasi-kebutuhan-sistem}

\section{Pengujian Kebutuhan Fungsional}

Tabel berikut menyajikan rincian pengujian kebutuhan fungsional pada alur Rawat Jalan dan Rawat Inap.

\begingroup
\small
\renewcommand{\arraystretch}{1.2}
\begin{longtable}{|p{0.12\textwidth}|p{0.32\textwidth}|p{0.28\textwidth}|p{0.12\textwidth}|}
	\caption{Pengujian Kebutuhan Fungsional pada Alur Rawat Jalan dan Rawat Inap}
	\label{tab:pengujian-kebutuhan-fungsional}                                                                                                 \\
	\hline
	\textbf{ID}                                                                        & \textbf{Skenario} & \textbf{Ekspektasi} & \textbf{KF} \\ \hline
	\endfirsthead

	\caption[]{Pengujian Kebutuhan Fungsional pada Alur Rawat Jalan dan Rawat Inap (lanjutan)}                                                 \\
	\hline
	\textbf{ID}                                                                        & \textbf{Skenario} & \textbf{Ekspektasi} & \textbf{KF} \\ \hline
	\endhead

	\texttt{P-KF-01}                                                                   &
	Pasien memberikan akses awal kepada petugas RM sesuai konteks episode pelayanan.   &
	Token awal episode pelayanan berhasil diterbitkan.                                 &
	\texttt{KF-TA-01}                                                                                                                          \\ \hline
	\texttt{P-KF-02}                                                                   &
	Petugas RM mendelegasikan akses pemeriksaan awal kepada perawat.                   &
	Token turunan perawat berhasil dibuat.                                             &
	\texttt{KF-TA-04}                                                                                                                          \\ \hline

	\texttt{P-KF-03}                                                                   &
	Petugas RM mendelegasikan akses klinis kepada dokter.                              &
	Token turunan dokter berhasil dibuat.                                              &
	\texttt{KF-TA-04}                                                                                                                          \\ \hline

	\texttt{P-KF-04}                                                                   &
	Perawat mengisi hasil pemeriksaan awal.                                            &
	Data pemeriksaan awal berhasil disimpan.                                           &
	\texttt{KF-TA-07}                                                                                                                          \\ \hline

	\texttt{P-KF-05}                                                                   &
	Dokter membaca pemeriksaan awal dan mengisi data klinis.                           &
	Data klinis berhasil dibaca dan disimpan.                                          &
	\texttt{KF-TA-05}                                                                                                                          \\ \hline

	\texttt{P-KF-06}                                                                   &
	Dokter mendelegasikan akses laboratorium.                                          &
	Token turunan laboratorium berhasil dibuat.                                        &
	\texttt{KF-TA-06}                                                                                                                          \\ \hline

	\texttt{P-KF-07}                                                                   &
	Petugas laboratorium membaca pemeriksaan penunjang dan mengisi hasil laboratorium. &
	Data pemeriksaan penunjang berhasil dibaca dan hasil laboratorium disimpan.        &
	\texttt{KF-TA-08}                                                                                                                          \\ \hline

	\texttt{P-KF-08}                                                                   &
	Dokter membaca hasil laboratorium.                                                 &
	Hasil laboratorium berhasil dibaca oleh dokter.                                    &
	\texttt{KF-TA-05}                                                                                                                          \\ \hline

	\texttt{P-KF-09}                                                                   &
	Dokter mendelegasikan akses apotek.                                                &
	Token turunan apoteker berhasil dibuat.                                            &
	\texttt{KF-TA-06}                                                                                                                          \\ \hline

	\texttt{P-KF-10}                                                                   &
	Apoteker membaca resep dan mengisi status pelayanan obat.                          &
	Data resep dan pelayanan obat berhasil diproses.                                   &
	\texttt{KF-TA-09}                                                                                                                          \\ \hline

	\texttt{P-KF-11}                                                                   &
	Pasien membuka riwayat delegasi episode pelayanan.                                 &
	Riwayat delegasi episode pelayanan ditampilkan.                                    &
	\texttt{KF-TA-03}                                                                                                                          \\ \hline
	\texttt{P-KF-12}                                                                   &
	Pasien merevokasi salah satu akses delegasi episode pelayanan.                     &
	Aktor \textit{delegatee} tidak bisa mengakses data RME pasien lagi.                &
	\texttt{KF-TA-02}                                                                                                                          \\ \hline
\end{longtable}
\endgroup

\section{Evaluasi Kebutuhan Sistem}

Tabel berikut menyajikan pemetaan kebutuhan sistem terhadap implementasi dan pengujian.

\begingroup
\small
\setlength{\tabcolsep}{3pt}
\begin{longtable}{|p{0.14\textwidth}|p{0.32\textwidth}|p{0.30\textwidth}|p{0.14\textwidth}|}
	\caption{Evaluasi Kebutuhan Sistem terhadap Implementasi dan Pengujian}
	\label{tab:evaluasi-kebutuhan-sistem}                                                                                                                                    \\
	\hline
	\textbf{Kebutuhan} & \textbf{Implementasi Terkait}                        & \textbf{Pengujian Terkait}                                                 & \textbf{Status} \\ \hline
	\endfirsthead
	\caption[]{Evaluasi Kebutuhan Sistem terhadap Implementasi dan Pengujian (lanjutan)}                                                                                     \\
	\hline
	\textbf{Kebutuhan} & \textbf{Implementasi Terkait}                        & \textbf{Pengujian Terkait}                                                 & \textbf{Status} \\ \hline
	\endhead
	\texttt{KF-TA-01}  & \texttt{I-TA-02}, \texttt{I-TA-03}, \texttt{I-TA-05} & \texttt{P-KF-01}                                                           & Terpenuhi.      \\ \hline
	\texttt{KF-TA-02}  & \texttt{I-TA-03}, \texttt{I-TA-04}, \texttt{I-TA-05} & \texttt{P-KF-12}, \texttt{P-KNF-11}, \texttt{P-KNF-12}                     & Terpenuhi.      \\ \hline
	\texttt{KF-TA-03}  & \texttt{I-TA-04}, \texttt{I-TA-05}                   & \texttt{P-KF-11}                                                           & Terpenuhi.      \\ \hline
	\texttt{KF-TA-04}  & \texttt{I-TA-02}, \texttt{I-TA-04}, \texttt{I-TA-06} & \texttt{P-KF-02}, \texttt{P-KF-03}                                         & Terpenuhi.      \\ \hline
	\texttt{KF-TA-05}  & \texttt{I-TA-01}, \texttt{I-TA-03}, \texttt{I-TA-06} & \texttt{P-KF-05}, \texttt{P-KF-08}                                         & Terpenuhi.      \\ \hline
	\texttt{KF-TA-06}  & \texttt{I-TA-02}, \texttt{I-TA-04}, \texttt{I-TA-06} & \texttt{P-KF-06}, \texttt{P-KF-09}                                         & Terpenuhi.      \\ \hline
	\texttt{KF-TA-07}  & \texttt{I-TA-01}, \texttt{I-TA-03}, \texttt{I-TA-06} & \texttt{P-KF-04}                                                           & Terpenuhi.      \\ \hline
	\texttt{KF-TA-08}  & \texttt{I-TA-01}, \texttt{I-TA-03}, \texttt{I-TA-06} & \texttt{P-KF-07}                                                           & Terpenuhi.      \\ \hline
	\texttt{KF-TA-09}  & \texttt{I-TA-01}, \texttt{I-TA-03}, \texttt{I-TA-06} & \texttt{P-KF-10}                                                           & Terpenuhi.      \\ \hline
	\texttt{KNF-TA-01} & \texttt{I-TA-01}, \texttt{I-TA-02}, \texttt{I-TA-03} & \texttt{P-KNF-01}, \texttt{P-KNF-02}                                       & Terpenuhi.      \\ \hline
	\texttt{KNF-TA-02} & \texttt{I-TA-02}, \texttt{I-TA-03}, \texttt{I-TA-06} & \texttt{P-KNF-03}, \texttt{P-KNF-04}, \texttt{P-KNF-05}, \texttt{P-KNF-06} & Terpenuhi.      \\ \hline
	\texttt{KNF-TA-03} & \texttt{I-TA-02}, \texttt{I-TA-03}, \texttt{I-TA-06} & \texttt{P-KNF-07}, \texttt{P-KNF-08}, \texttt{P-KNF-09}                    & Terpenuhi.      \\ \hline
	\texttt{KNF-TA-04} & \texttt{I-TA-04}, \texttt{I-TA-05}, \texttt{I-TA-06} & \texttt{P-KF-11}, \texttt{P-KNF-10}                                        & Terpenuhi.      \\ \hline
	\texttt{KNF-TA-05} & \texttt{I-TA-03}, \texttt{I-TA-05}                   & \texttt{P-KNF-11}, \texttt{P-KF-12}, \texttt{P-KNF-12}                     & Terpenuhi.      \\ \hline
	\texttt{KNF-TA-06} & \texttt{I-TA-03}, \texttt{I-TA-06}                   & \texttt{P-KNF-09}, \texttt{P-KNF-13}                                       & Terpenuhi.      \\ \hline
\end{longtable}
\endgroup
